% ============================================================================
% LEHRERHINWEISE - Physik Klasse 6: Weg-Zeit-Diagramme
% ============================================================================
% Kurz-Lessonplan + Lösungen + Differenzierungshinweise
% ============================================================================

\documentclass[11pt,a4paper]{article}

\usepackage[utf8]{inputenc}
\usepackage[T1]{fontenc}
\usepackage[ngerman]{babel}
\usepackage{geometry}
\usepackage{helvet}
\usepackage{tikz}
\usepackage{amsmath}
\usepackage{enumitem}
\usepackage{tabularx}
\usepackage{array}
\usepackage{longtable}
\usepackage{fancyhdr}
\usepackage{lastpage}
\usepackage{xcolor}
\usepackage{tcolorbox}

\renewcommand{\familydefault}{\sfdefault}

\geometry{left=2cm, right=2cm, top=2.5cm, bottom=2cm}

% Farben
\definecolor{physik}{HTML}{2c5aa0}
\definecolor{success}{HTML}{2a7a4b}
\definecolor{warning}{HTML}{b35c00}
\definecolor{error}{HTML}{c62828}
\definecolor{lightgray}{HTML}{f5f5f5}

\tcbuselibrary{skins,breakable}

\newtcolorbox{tippbox}{
    colback=lightgray,
    colframe=success,
    fonttitle=\bfseries,
    title=Tipp,
    breakable
}

\newtcolorbox{warnbox}{
    colback=lightgray,
    colframe=warning,
    fonttitle=\bfseries,
    title=Achtung,
    breakable
}

\newtcolorbox{fehlerbox}{
    colback=lightgray,
    colframe=error,
    fonttitle=\bfseries,
    title=Häufige Fehler,
    breakable
}

% Kopfzeile
\pagestyle{fancy}
\fancyhf{}
\fancyhead[L]{\small\textbf{LEHRERHINWEISE} | Physik Klasse 6 | Weg-Zeit-Diagramme}
\fancyhead[R]{\small Stunde 3 von 4}
\fancyfoot[C]{\small Seite \thepage{} von \pageref{LastPage}}
\renewcommand{\headrulewidth}{0.4pt}

\begin{document}

% ============================================================================
% TITEL
% ============================================================================
\begin{center}
    {\Large\bfseries\textcolor{physik}{Lehrerhinweise: Weg-Zeit-Diagramme}}\\[0.3em]
    {\normalsize Physik Klasse 6 | Stunde 3 von 4 | ca. 45 Minuten}
\end{center}

\vspace{0.5em}
\hrule
\vspace{1em}

% ============================================================================
% ÜBERSICHT
% ============================================================================
\section*{Übersicht}

\begin{tabular}{ll}
    \textbf{Thema:} & Weg-Zeit-Diagramme verstehen und zeichnen \\
    \textbf{Reihe:} & Bewegungen -- Geradlinige, gleichförmige Bewegung \\
    \textbf{Vorwissen:} & Bewegungsarten (Std. 1), Geschwindigkeit $v = s/t$ (Std. 2) \\
    \textbf{Materialien:} & AB-003, LP-003 (optional), Lineal, Tafel/Whiteboard \\
    \textbf{Rahmenplan:} & MV Kl. 6 OS: ``s(t)-Diagramme -- Graphische Darstellung möglich''
\end{tabular}

\section*{Lernziele}

\begin{itemize}[leftmargin=2em]
    \item Die SuS können die Achsen eines s(t)-Diagramms korrekt beschriften.
    \item Die SuS können aus einem s(t)-Diagramm Werte für Weg und Zeit ablesen.
    \item Die SuS können aus einer Wertetabelle ein s(t)-Diagramm zeichnen.
    \item Die SuS können qualitativ erklären, was eine steile bzw. flache Gerade bedeutet.
\end{itemize}

% ============================================================================
% STUNDENVERLAUF
% ============================================================================
\section*{Stundenverlauf (45 Minuten)}

\begin{longtable}{|p{1cm}|p{1.5cm}|p{6cm}|p{2.5cm}|p{2cm}|}
\hline
\textbf{Zeit} & \textbf{Phase} & \textbf{Inhalt} & \textbf{Sozialform} & \textbf{Material} \\
\hline
\endhead

0--5 & Einstieg & L zeigt Wertetabelle aus Stunde 2. Impuls: ``Wie können wir das übersichtlicher darstellen?'' SuS nennen Ideen (Diagramm). & Plenum & Tafel \\
\hline

5--10 & Input I & L erklärt Grundaufbau s(t)-Diagramm: Achsen beschriften (t waagerecht, s senkrecht), Einheiten. SuS füllen Kasten 1 im AB aus. & LSG & Tafel, AB \\
\hline

10--18 & Input II & L demonstriert: Wertetabelle $\to$ Punkte einzeichnen $\to$ Gerade verbinden. Bespricht Interpretation (steil = schnell). SuS füllen Kasten 2 aus. & LSG & Tafel, AB \\
\hline

18--30 & Übung & SuS bearbeiten Aufgaben 1--3 (differenziert). L geht herum, gibt individuelle Hilfe. & EA & AB, Lineal \\
\hline

30--35 & Partner-check & SuS vergleichen Ergebnisse mit Nachbar. Diskutieren Unterschiede. & PA & AB \\
\hline

35--42 & Sicherung & L bespricht Lösungen (Doku-Kamera/Tafel). Klärt typische Fehler. Wiederholt Merksätze. & Plenum & Doku-Kamera \\
\hline

42--45 & Abschluss & Ausblick auf Stundenleistung. HA: Aufgabe 4 oder Lernpfad (optional). & Plenum & -- \\
\hline
\end{longtable}

% ============================================================================
% DIFFERENZIERUNG
% ============================================================================
\section*{Differenzierung}

\begin{tabular}{|l|p{4cm}|p{4cm}|p{4cm}|}
\hline
& \textbf{Basis (*)} & \textbf{Standard (**)} & \textbf{Erweitert (***)} \\
\hline
\textbf{Fokus} & Achsen beschriften, Werte ablesen & Diagramm zeichnen, interpretieren & Vergleichen, Geschichte erfinden \\
\hline
\textbf{AB} & Kästen 1--2, Aufgabe 1 & + Aufgabe 2, 3 & + Aufgabe 4 \\
\hline
\textbf{Hilfen} & Koordinatensystem vorgegeben, Punkte teilweise markiert & Hilfslinien im Koordinatensystem & Keine Hilfen \\
\hline
\textbf{Erwartung} & 10 von 10 P erreichbar & 20+ von 24 P angestrebt & 25+ von 29 P angestrebt \\
\hline
\end{tabular}

\vspace{1em}

\begin{tippbox}
\textbf{Für leistungsschwache SuS:}
\begin{itemize}[leftmargin=1.5em]
    \item Koordinatensystem mit eingezeichneten Hilfslinien bereitstellen
    \item Ersten Punkt gemeinsam einzeichnen
    \item Wortliste: ``waagerecht'', ``senkrecht'', ``steil'', ``flach''
\end{itemize}
\end{tippbox}

\begin{tippbox}
\textbf{Für leistungsstarke SuS:}
\begin{itemize}[leftmargin=1.5em]
    \item Aufgabe 4 als Pflicht (nicht optional)
    \item Zusatz: Zwei Bewegungen ins gleiche Diagramm zeichnen
    \item Expertenrolle: Anderen SuS beim Zeichnen helfen
\end{itemize}
\end{tippbox}

% ============================================================================
% HÄUFIGE FEHLER
% ============================================================================
\section*{Häufige Fehler und Reaktionen}

\begin{fehlerbox}
\textbf{1. Achsen vertauscht (t auf y-Achse)}

\textit{Reaktion:} ``Merke dir: Die Zeit läuft von links nach rechts, wie beim Lesen. Deshalb kommt sie auf die waagerechte Achse.''

\vspace{0.5em}
\textbf{2. Graph als Abbild der Bewegung (``Der fährt bergauf'')}

\textit{Reaktion:} ``Das Diagramm zeigt nicht die Landschaft! Eine steigende Gerade bedeutet, dass der Weg zunimmt, nicht dass jemand bergauf fährt.''

\vspace{0.5em}
\textbf{3. Falsche Skalierung (ungleichmäßige Abstände)}

\textit{Reaktion:} ``Achte darauf, dass die Abstände auf der Achse immer gleich sind. Von 0 zu 10 und von 10 zu 20 muss der gleiche Abstand sein.''

\vspace{0.5em}
\textbf{4. Punkte nicht auf Gitternetz (ungenaues Einzeichnen)}

\textit{Reaktion:} ``Lies zuerst den t-Wert auf der x-Achse, dann gehe senkrecht hoch. Lies dann den s-Wert und gehe waagerecht. Wo sich beide treffen, ist dein Punkt.''

\vspace{0.5em}
\textbf{5. Verwechslung steil/flach $\leftrightarrow$ schnell/langsam}

\textit{Reaktion:} ``Bei einer steilen Geraden legt man in kurzer Zeit viel Weg zurück. Das ist eine schnelle Bewegung. Probier es selbst: In 2 Sekunden 20 Meter -- schnell oder langsam?''
\end{fehlerbox}

% ============================================================================
% LÖSUNGEN
% ============================================================================
\section*{Lösungen}

\subsection*{Kasten 1: Achsenbeschriftung}
\begin{itemize}[leftmargin=2em]
    \item Waagerechte Achse (x-Achse): \textbf{Zeit $t$} in \textbf{Sekunden (s)}
    \item Senkrechte Achse (y-Achse): \textbf{Weg $s$} in \textbf{Metern (m)}
\end{itemize}

\subsection*{Kasten 2: Interpretation}
\begin{itemize}[leftmargin=2em]
    \item Steile Gerade = \textbf{schnelle} Bewegung
    \item Flache Gerade = \textbf{langsame} Bewegung
    \item Waagerechte Linie = \textbf{Stillstand}
\end{itemize}

\subsection*{Aufgabe 1 (6 Punkte)}
\begin{enumerate}[label=\alph*)]
    \item Nach $t = 4$ s: $s = \textbf{40 m}$ \hfill (2 P)
    \item Nach $t = 10$ s: $s = \textbf{100 m}$ \hfill (2 P)
    \item Für $s = 60$ m: $t = \textbf{6 s}$ \hfill (2 P)
\end{enumerate}

\subsection*{Aufgabe 2 (8 Punkte)}
\begin{itemize}[leftmargin=2em]
    \item a) + b) Diagramm: 5 Punkte korrekt eingezeichnet (5 P), Gerade verbunden (1 P)
    \item c) Geschwindigkeit: $v = \dfrac{100\,\text{m}}{20\,\text{s}} = \textbf{5 m/s}$ \hfill (2 P)
\end{itemize}

\begin{center}
\begin{tikzpicture}[scale=0.5]
    \draw[->] (0,0) -- (11,0) node[right] {\small $t$ in s};
    \draw[->] (0,0) -- (0,6) node[above] {\small $s$ in m};
    \foreach \x/\label in {0/0, 2.5/5, 5/10, 7.5/15, 10/20} {
        \draw (\x,-0.1) -- (\x,0.1);
        \node[below] at (\x,-0.3) {\tiny \label};
    }
    \foreach \y/\label in {0/0, 1.25/25, 2.5/50, 3.75/75, 5/100} {
        \draw (-0.1,\y) -- (0.1,\y);
        \node[left] at (-0.3,\y) {\tiny \label};
    }
    \draw[thick, physik] (0,0) -- (10,5);
    \fill[physik] (0,0) circle (3pt);
    \fill[physik] (2.5,1.25) circle (3pt);
    \fill[physik] (5,2.5) circle (3pt);
    \fill[physik] (7.5,3.75) circle (3pt);
    \fill[physik] (10,5) circle (3pt);
\end{tikzpicture}
\end{center}

\subsection*{Aufgabe 3 (6 Punkte)}
\begin{enumerate}[label=\alph*)]
    \item Schneller: \textbf{Person A} \hfill (1 P)
    \item Begründung: \textbf{Die Gerade von A ist steiler als die von B. Eine steilere Gerade bedeutet eine höhere Geschwindigkeit.} \hfill (3 P)
    \item Nach 8 s: $s = \textbf{80 m}$ \hfill (2 P)
\end{enumerate}

\subsection*{Aufgabe 4 (5 Punkte)}
\textbf{Erwartungshorizont:}
\begin{itemize}[leftmargin=2em]
    \item Geschichte passt zum Diagrammverlauf (3 P)
    \item Drei Phasen erkannt (schnell -- Stillstand -- langsam) (2 P)
\end{itemize}

\textbf{Beispiellösung:}

``Lisa fährt mit dem Fahrrad zur Schule. In den ersten 3 Minuten fährt sie schnell bergab und legt 300 m zurück. Dann muss sie an einer roten Ampel warten (3 Minuten Stillstand). Danach fährt sie langsamer weiter, weil es bergauf geht, und erreicht nach weiteren 4 Minuten die Schule bei 500 m.''

% ============================================================================
% MATERIAL
% ============================================================================
\section*{Benötigtes Material}

\begin{itemize}[leftmargin=2em]
    \item Arbeitsblatt AB-003 (Klassensatz)
    \item Lernpfad LP-003 (optional, für Tablets/PCs oder Hausaufgabe)
    \item Lineal (für jeden Schüler)
    \item Tafel/Whiteboard mit Koordinatensystem
    \item Dokumentenkamera (für Besprechung)
    \item Farbige Kreide/Stifte (zur Unterscheidung verschiedener Bewegungen)
\end{itemize}

% ============================================================================
% AUSBLICK
% ============================================================================
\section*{Ausblick: Stunde 4 (Stundenleistung)}

Die nächste Stunde dient der Anwendung und Überprüfung:
\begin{itemize}[leftmargin=2em]
    \item Stundenleistung (ca. 30 Minuten)
    \item Differenziert in drei Niveaus (*/~**/~***)
    \item Aufgabentypen: Klassifikation, Ablesen, Zeichnen, Vergleichen
    \item Punkte: (*) ca. 26 P, (**) ca. 31 P, (***) ca. 35 P
\end{itemize}

\begin{warnbox}
\textbf{Vorbereitung für Stunde 4:}
\begin{itemize}[leftmargin=1.5em]
    \item Stundenleistungsbögen kopieren (3 Versionen)
    \item Ggf. Nachteilsausgleich für LRS-SuS berücksichtigen (Zeitzugabe)
    \item Lineal-Verfügbarkeit sicherstellen
\end{itemize}
\end{warnbox}

\end{document}
